% noman_report.tex — No-Man causal risk report LaTeX template
% Requires: XeLaTeX, tcolorbox (most), tikz, xeCJK, fontspec, enumitem
% Variables are substituted via \VAR{name} string replacement in noman_latex.py.
\documentclass[a4paper,11pt]{article}
\usepackage{fontspec}
\usepackage{xeCJK}
\setCJKmainfont{Hiragino Mincho ProN}
\setCJKsansfont{Hiragino Sans}
\usepackage[a4paper, top=2cm, bottom=2cm, left=2.5cm, right=2.5cm]{geometry}
\usepackage[most]{tcolorbox}
\usepackage{tikz}
\usetikzlibrary{positioning, arrows.meta}
\usepackage{xcolor}
\usepackage{parskip}
\usepackage{fancyhdr}
\usepackage{enumitem}
\usepackage{hyperref}
\setlength{\parindent}{0pt}
\setlength{\parskip}{4pt}

% ── Confidence tier colours ────────────────────────────────────────────────────
\definecolor{tierwell}{RGB}{34,139,34}
\definecolor{tierplausible}{RGB}{30,100,200}
\definecolor{tierworth}{RGB}{190,110,0}
\definecolor{tiernone}{RGB}{120,120,120}

% ── Severity badge colours ─────────────────────────────────────────────────────
\definecolor{sevhigh}{RGB}{180,30,30}
\definecolor{sevmedium}{RGB}{190,110,0}
\definecolor{sevlow}{RGB}{60,130,60}

\pagestyle{fancy}
\fancyhf{}
\rhead{\small No-Man 因果リスク分析}
\lhead{\small \VAR{decision_label}}
\cfoot{\thepage}

\begin{document}

% ── Dashboard strip ───────────────────────────────────────────────────────────
\begin{tcolorbox}[
  enhanced,
  colback=gray!10,
  colframe=gray!55,
  boxrule=0.8pt,
  arc=4pt,
  left=10pt, right=10pt, top=8pt, bottom=8pt,
  fonttitle=\bfseries\large,
  title={No-Man\quad 因果リスク分析レポート},
]
\renewcommand{\arraystretch}{1.35}
\begin{tabular}{@{}p{5.5cm}p{8cm}@{}}
  \textbf{意思決定タイプ}       & \VAR{decision_label}  \\
  \textbf{生成日}               & \VAR{report_date}     \\
  \textbf{検出連鎖数}           & \VAR{chain_count}     \\
  \textbf{グラフバージョン}     & \VAR{graph_version}   \\
  \textbf{プロモーションモード} & \VAR{promotion_mode}  \\
\end{tabular}
\end{tcolorbox}

\vspace{0.7em}

% ── Executive summary ─────────────────────────────────────────────────────────
\begin{tcolorbox}[
  enhanced,
  colback=gray!4,
  colframe=gray!40,
  boxrule=0.5pt,
  arc=4pt,
  fonttitle=\bfseries,
  title={経営陣向けサマリー},
]
\VAR{executive_summary}
\end{tcolorbox}

\vspace{0.7em}

% ── Causal chain DAG ─────────────────────────────────────────────────────────
\begin{tcolorbox}[
  enhanced,
  colback=white,
  colframe=gray!40,
  boxrule=0.5pt,
  arc=4pt,
  fonttitle=\bfseries,
  title={因果連鎖グラフ（全連鎖）},
]
\begin{center}
\resizebox{\linewidth}{!}{%
\VAR{tikz_dag}
}
\end{center}
\begin{center}
{\small \textcolor{blue}{$\rightarrow$} 正の効果\quad \textcolor{red}{$\rightarrow$} 負の効果}
\end{center}
\end{tcolorbox}

\clearpage

% ── Chain detail blocks ───────────────────────────────────────────────────────
\section*{詳細分析}

\VAR{chain_blocks}

\end{document}
